% !TEX root = ../main.tex

\chapter{基于分层架构与 MuJoCo 仿真的四足机器人视觉导航闭环系统设计与验证}

\section{问题分析}

第二章已经给出了可进入闭环验证的高层视觉导航算法配置，但算法本身并不等同于机器人系统。NoMaD 输出的是局部航点轨迹，而真实或仿真的四足机器人需要的是连续、限幅、可安全中断的运动指令；离线实验中的一次模型前向，也无法直接处理任务切换、目标节点更新、异常恢复和平台通信等系统问题。因此，本章面对的核心问题是：在已有动作扩散视觉导航算法的基础上，如何构建一个完整的四足机器人视觉导航闭环系统，并在进入真机之前验证它是否能够稳定运行。

这一问题包含三个层面的矛盾。首先是算法输出与平台执行之间的矛盾。高层策略给出的航点具有平台无关性，但 Lite3 运动主机或 MuJoCo 后端最终接收的是速度参考，二者之间必须有明确的中层控制映射和限幅机制。其次是单次推理与连续任务之间的矛盾。视觉导航不是孤立调用模型，而是伴随拓扑地图节点选择、导航/探索模式切换、完成判定、恢复和急停等状态变化，需要状态机和导航主机把这些逻辑组织起来。最后是仿真验证与真实部署之间的矛盾。直接上真机会受到场地、安全和硬件风险限制，因此需要先在 MuJoCo 中以可控方式检查接口、时序和闭环行为。

基于上述分析，本章将第二章得到的高层推理配置封装为统一推理模块，并在其外部构建状态机、导航主机、中层控制映射和平台接口。MuJoCo 后端作为该平台接口的第一个具体落地对象，用于检验从视觉观察、目标选择、轨迹生成、速度映射到平台执行的完整链路。通过这种组织方式，本章不只是给出仿真实验结果，而是完成从算法层到系统层、再到仿真验证层的过渡，为第四章 Jetson Orin 与 Lite3 真机部署提供可复用接口、初始参数和安全边界。

\section{四足机器人平台与闭环接口基础}

\subsection{轮式与足式平台的部署差异}

轮式平台姿态变化小、视觉输入稳定，与速度控制器的接口也较简单；足式平台则具有更强的地形适应性，但会带来更明显的机体起伏、更严格的实时性约束和更复杂的控制接口\cite{hwangbo2019anymal,lee2020learning}。这意味着，同一高层轨迹若能在轮式平台上顺利执行，并不代表它能直接在四足平台上稳定工作。

具体而言，足式平台的步态运动会在每个步态周期内引入机体俯仰和侧倾变化，这种扰动会导致相机的视场发生抖动，使得连续帧之间的视觉特征产生非平移性变化。对于依赖时序视觉输入的导航策略，步态引起的视觉抖动可能被误解读为环境变化或自身运动，从而产生错误的轨迹预测。此外，足式平台的非完整运动学约束也与轮式平台不同。轮式机器人通常可以精确执行差速运动学模型所描述的运动，而四足机器人的实际运动轨迹受步态规划器和地形交互的影响，与指令速度之间存在非线性偏差。这种偏差意味着即使高层策略输出了精确的速度指令，机器人的实际位移也可能与预期存在差异，进一步强调了闭环执行和频繁重规划的必要性。


\subsection{云深处 Lite3 平台与关键参数}

本文采用的主要足式平台为云深处 Lite3。结合本文系统设计与部署需求，
其在本文中的系统定位是``由高层视觉导航策略驱动的四足运动系统''，而不是端到端关节控制平台。对应控制链路为
\begin{equation}
\text{视觉观察} \rightarrow \text{高层 NoMaD 推理} \rightarrow \text{中层控制映射} \rightarrow \text{底层运动执行}.
\end{equation}
式中，箭头表示信息流和控制流的传递方向；“视觉观察”对应相机图像输入，“高层 NoMaD 推理”对应局部航点生成，“中层控制映射”对应速度参考计算，“底层运动执行”对应 Lite3 运动主机或仿真后端的实际执行。

与本文系统设计和实验验证相关的 Lite3 关键参数如表~\ref{tab:lite3-params} 所示。
\begin{table}[htbp]
    \centering
    \caption{本文中与 Lite3 相关的关键平台参数}
    \label{tab:lite3-params}
    \begin{tabular}{lll}
        \toprule
        参数类别 & 参数项 & 数值或说明 \\
        \midrule
        平台类型 & 机器人形态 & 四足平台 \\
        运动结构 & 关节数量 & 12 关节执行系统 \\
        速度接口 & 前后速度范围 & $[-1.0,1.0]\,\text{m/s}$ \\
        速度接口 & 左右速度范围 & $[-0.5,0.5]\,\text{m/s}$ \\
        速度接口 & 偏航角速度范围 & $[-1.5,1.5]\,\text{rad/s}$ \\
        部署限幅 & 本文最大线速度 & $0.12\,\text{m/s}$ \\
        部署限幅 & 本文最大角速度 & $0.35\,\text{rad/s}$ \\
        高层闭环 & 仿真 NoMaD 周期 & $4\,\text{Hz}$ \\
        通信安全 & 速度超时保护 & $250\,\text{ms}$ 自动停止 \\
        \bottomrule
    \end{tabular}
\end{table}


\begin{figure}[htbp]
    \centering
    \includegraphics[width=0.5\linewidth]{prelim-lite3-placeholder.png}
    \caption{Lite3 平台与控制链路示意图}
    \label{fig:prelim-lite3-placeholder}
\end{figure}

在通信协议层面，本文系统与 Lite3 之间通过 UDP 进行数据交换。桥接层将高层三维速度参考转换为前后、左右和偏航三个速度通道，并持续维护心跳与速度刷新。本文设置心跳频率为 $5\,\text{Hz}$，速度刷新频率为 $25\,\text{Hz}$，满足 Lite3 通讯文档中速度轴命令至少 $20\,\text{Hz}$ 刷新的要求。协议文档给出的速度命令超时约束可以理解为底层接口的安全边界，但高层 NoMaD 推理本身并不需要在每个 $250\,\text{ms}$ 内都完整完成一次；桥接层会在两次高层推理之间持续重复最近一次限幅后的速度参考。

关于观察频率与控制频率的关系，需要区分三个层次的频率，这种多层频率架构可以用以下层级关系概括：
\begin{equation}
f_{\text{joint}} \gg f_{\text{velocity}} \gg f_{\text{vision}},
\end{equation}
其中 $f_{\text{joint}}$ 表示底层关节级控制频率，$f_{\text{velocity}}$ 表示速度命令刷新频率，本文构建的真机桥接层中 $f_{\text{velocity}}=25\,\text{Hz}$，$f_{\text{vision}}$ 表示高层视觉导航策略的更新频率，在 MuJoCo 仿真中为 $4\,\text{Hz}$，在 Orin 真机上则由实测推理耗时决定。这种频率分离设计的优势是每一层只需关注自身时间尺度上的决策，代价是层间的响应延迟会累积。

\section{分层闭环系统架构设计}\label{sec:unified-arch}

\subsection{分层设计}

本文将系统划分为四层：推理层、任务组织层、中层控制层和平台执行层。推理层负责根据视觉观察与目标图像生成局部轨迹；任务组织层负责状态切换、任务队列与模式切换；中层控制层负责将轨迹映射为速度或姿态参考；平台执行层负责向 MuJoCo 或真机后端下发指令。

这种四层划分的设计动机源自足式导航系统的本质需求。推理层的存在使得高层策略可以独立于具体平台进行开发和优化：无论底层是四足机器人还是轮式平台，高层 NoMaD 生成的局部轨迹格式保持一致。
任务组织层的引入是因为实际部署中，导航任务往往不是单次调用，而是涉及多目标队列、模式切换和异常恢复等复杂流程，这些逻辑不应侵入推理层或控制层。
中层控制层的核心价值在于不同平台的可复用性。不同的机器人本体或者是导航平台，运动控制的实现可能是不同的，但是接收速度指令的机制是相同的。
平台执行层则直接面向硬件或仿真后端，负责指令下发和传感器回传。

与单体式设计相比，分层架构虽然引入了层间通信开销，但带来的解耦优势远大于性能损失。在单体式设计中，推理逻辑、任务组织和平台控制混杂在同一模块内，任何局部修改都可能引发全局回归。
例如，当需要将 DDIM 步数从 4 调整为 2 时，单体式系统可能需要同时修改推理、控制和部署流程，而在分层架构中，这一变更完全限制在推理层内部。
此外，单体式设计在面对多平台部署时需要维护多套几乎重复的实现，而分层架构只需替换平台执行层即可。从工程实践角度来看，分层架构还简化了调试流程：当闭环出现问题时，可以通过逐层隔离快速定位故障所在层级。

层间接口方面，推理层向任务组织层输出的是轨迹候选、候选节点索引和距离预测值；任务组织层根据当前状态选择是否继续导航、探索、恢复或停机；中层控制层接收当前选定的航点，并输出统一的速度参考；平台执行层再将该参考映射到 MuJoCo 命令或真机桥接命令。

\begin{figure}[htbp]
    \centering
    \includegraphics[width=0.8\linewidth]{fenceng.png}
    \caption{分层闭环系统示意图}
    \label{fig:fenceng}
\end{figure}


\subsection{统一推理模块}

统一推理模块集中管理模型权重、采样器类型、采样步数、条件引导强度和候选轨迹预算等推理设置。
其作用是把 DDIM、CFG、TTS 和视觉编码器替换都限制在高层策略层，使状态机、导航主机与平台后端不需要感知采样器和编码器细节。该模块直接承接第二章实验最终确定的``EfficientNet-B0 + DDIM-2 + CFG=0 + TTS-8''默认配置；统一推理模块对外暴露的标准推理流程也与第二章算法~\ref{algo:inference-pipeline} 描述的``视觉编码 $\to$ DDIM 去噪 $\to$ CFG 外推 $\to$ TTS 评分''层级一致。这样，第二章的算法结论被完整地提炼为系统层的可复用模块，而不再以分散的实验流程形式存在。


参数变更在系统中的传播路径是单向且受限的。采样步数或条件引导强度的调整只影响推理模块内部的轨迹生成过程，不会传播到状态机的状态转移逻辑或平台后端的指令格式。推理模块输出的轨迹分布可能随参数变化而改变，但下游模块只关心最终选定的轨迹点，而不关心该轨迹的具体采样细节。
这种封装策略确保了高层算法的快速迭代不会引发系统级回归。在实验过程中，本文多次调整采样参数，均未需要修改状态机或平台后端逻辑，验证了这一设计的有效性。

这意味着：只要新的策略仍保持兼容接口，后续系统模块就可以复用。也正因如此，本文将「推理模块」单独提炼为系统组件，而不是把 DDIM、CFG 或 TTS 的逻辑分布于多个实验流程中。

\section{状态机与导航主机设计}\label{sec:fsm-host}

\subsection{状态机职责}

在足式闭环系统中，高层推理并不是单次调用，而是需要持续嵌入控制循环。因此，系统必须回答何时起立、何时导航、何时探索、何时恢复、何时停机等问题。
为此，本文完成了显式状态机设计。本文系统的核心状态包括空闲、起立、导航、探索、恢复、完成、失败和急停等。

状态机的作用不是增加流程复杂度，而是保证系统具有可恢复性和安全边界。对于 Lite3 这样的足式平台，一旦高层轨迹异常、控制映射失效或通信超时，系统必须进入恢复或停机状态，而不是继续盲目前进。

各状态的进入与退出条件围绕任务生命周期组织。空闲状态是每次任务内部的入口状态，随后进入起立状态完成平台准备；导航状态持续读取相机图像、更新上下文帧、调用 NoMaD 预测局部轨迹，并通过中层控制映射下发速度命令。在 MuJoCo 后端，若任务存在物理目标，则以机器人当前位置到目标点的距离小于 $0.5\,\mathrm{m}$ 作为到达判定；在真实拓扑地图导航中，则主要依赖拓扑地图节点选择和距离预测来判断是否接近目标节点。探索状态与导航状态共用推理链路，但通过 goal mask 构造无目标条件；恢复状态用于处理卡住等异常情况；急停状态则触发平台停机保护。


\begin{figure}[htbp]
    \centering
    \includegraphics[width=0.8\linewidth]{figframework-fsm-placeholder.png}
    \caption{状态机状态流转图}
    \label{fig:figframework-fsm-placeholder}
\end{figure}

\subsection{导航主机职责}

如果说状态机解决的是任务内部如何运行，那么导航主机解决的则是系统需要运行什么任务。导航主机负责选择仿真或真机后端、装载单目标或多目标任务、切换导航/探索模式，并将任务分发给状态机执行。

在任务组织方面，本文导航主机支持通过 JSON 格式配置批量读取任务，也支持实时控制入口。

在目标选择策略方面，对于单目标导航，系统加载一个与后端域匹配的拓扑地图数据路径，并在以当前定位指针为中心的局部搜索窗口内用距离预测头推进序列指针，再把序列上比该指针再前进一格的节点图像作为本周期送入 NoMaD 的目标条件；对于多目标仿真任务，系统可以在同一 MuJoCo 场景中采样多个目标点并为每个目标在线生成局部拓扑地图；对于真实部署，目标来源应是预先采集的真实拓扑地图数据。探索模式不提供真实目标图像，而是通过 goal mask 构造无目标条件，推理模块在该条件下生成探索轨迹。

模式切换逻辑方面，本文系统支持由操作者或任务配置显式切换不同任务；同一任务内部则由状态机在起立、导航或探索、恢复、完成或失败等状态之间切换。

这种分层使多目标导航不必由 NoMaD 一次性做全局规划，而是由导航主机把长任务拆成一系列局部子任务，再由高层策略逐段完成。由此，单目标导航、探索模式和多目标任务都能在同一高层策略之上被统一组织。

\subsection{中层控制映射}

中层控制映射是连接高层轨迹输出与平台执行指令的桥梁。
推理模块输出的是以机器人当前位姿为原点的局部轨迹点序列，通常包含未来若干步的二维位移预测。
中层控制映射模块从该序列中选取最近的轨迹点作为当前步的目标位移，并根据位移方向和幅度计算线速度与角速度指令。
对于 Lite3 四足平台，中层 PD 控制器首先把前向速度限制在 $[0,0.4]\,\mathrm{m/s}$，角速度限制在 $[-0.8,0.8]\,\mathrm{rad/s}$；真机部署时还会在中层提供线速度和角速度缩放参数，并由桥接层进一步采用更保守的安全限幅。本文真机配置为前向速度不超过 $0.12\,\mathrm{m/s}$、偏航角速度不超过 $0.35\,\mathrm{rad/s}$。
当轨迹点位移极小（低于 0.01m 阈值）时，中层映射输出零速度指令，避免因微小噪声导致的无意义运动。
这种映射策略简单但有效，将高层扩散策略的连续轨迹输出转化为离散的速度控制指令，同时通过速度限幅保证了平台安全性。

\section{MuJoCo 仿真平台与任务设计}\label{sec:mujoco-platform}

前文已经给出了四足机器人视觉导航闭环系统的模块边界，但系统设计本身还不足以说明链路能够稳定运行。对于真实 Lite3 平台而言，直接在真机上调试高层视觉策略、状态机和速度控制映射存在一定风险：错误的航点、过大的角速度或异常状态切换都可能导致碰撞、摔倒或通信失控。因此，本文先在 MuJoCo 中建立与分层架构兼容的仿真后端，使高层算法和中层控制能够在可重复、可观察的环境中完成闭环复核。

\subsection{选择 MuJoCo 的原因}

MuJoCo 具有较好的接触动力学建模能力和可调仿真接口，被广泛用于足式控制与学习系统验证\cite{mujoco2012,hwangbo2019anymal,rl_locomotion_rudin2022}。
对本文而言，MuJoCo 的价值不只是把机器人动起来，而是提供一个风险较低、接口清晰的闭环环境，用于验证高层推理、状态机、中层控制映射与平台后端是否能够协同工作。
此外，MuJoCo 支持无显示渲染模式（headless mode），可以在无 GPU 显示的服务器上批量运行仿真实验，这为后续的参数搜索和消融实验提供了便利。

\subsection{仿真系统组成与接口一致性}

本文采用的 Lite3 MuJoCo 仿真链路可以概括为五部分：场景与拓扑地图、NoMaD 推理模块、中层轨迹映射、MuJoCo 平台接口和状态机循环。
仿真后端为高层模块提供相机图像、姿态、位置信息和安全停机等能力，使其尽可能与真机后端保持一致。其中，MuJoCo 平台接口封装了仿真器的步进、渲染和状态读取功能，对上层暴露与真机后端一致的图像获取、状态读取、速度下发和急停等接口。
这种接口一致性是仿真到真机迁移的基础。

从本章的结构来看，MuJoCo 后端并不是脱离系统架构的独立实验，而是平台执行层的一种具体实现。统一推理模块仍然输出局部航点，状态机仍然负责起立、导航、探索、恢复和急停等状态切换，中层控制映射仍然把局部轨迹转化为速度参考；区别仅在于平台执行层把速度参考传递给 MuJoCo 中的 Lite3 模型，而不是真机运动主机。这样可以在不改动高层逻辑的情况下，检查分层接口是否真正具备双后端复用能力。

\subsection{场景、物理参数与相机设置}

仿真场景的设计需要在复杂度与可控性之间取得平衡。本文仿真中的三类场景均采用装饰走廊形式，并通过走廊长度、宽度和障碍物数量形成难度梯度。easy 场景为半宽 $1.2\,\mathrm{m}$、目标点约在 $x=5\,\mathrm{m}$ 的无障碍走廊；medium 场景为半宽 $1.5\,\mathrm{m}$、目标点约在 $x=7\,\mathrm{m}$ 的单障碍走廊；hard 场景在更长走廊中布置两个交错障碍，目标点约在 $x=8\,\mathrm{m}$。
每个场景都可生成对应的 MuJoCo 域拓扑地图节点图像，为导航任务提供目标参考。

物理与控制节拍方面，本文仿真设置 MuJoCo 物理步长为 $0.001\,\mathrm{s}$，强化学习（Reinforcement Learning, RL）运动策略周期为 $0.02\,\mathrm{s}$，即底层步态策略约 $50\,\mathrm{Hz}$ 更新一次；高层 NoMaD 控制周期为 $4\,\mathrm{Hz}$，每个 NoMaD 周期包含约 12 个 RL 策略步。

相机仿真方面，MuJoCo 内置渲染器输出第一人称 RGB 图像。本文仿真中的离屏渲染尺寸为 $320 \times 240$，随后由统一推理模块按模型配置缩放到 $96\times96$ 或其他编码器所需尺寸。本文没有把仿真相机严格标定到某一款 RGB 相机或者 RealSense 相机的视场角，因此将其定位为与真机前视相机功能相近的第一人称视角（First-Person View, FPV）观测，而非真实相机参数级对齐的仿真模型。需要注意的是，仿真相机不包含真实相机的径向畸变、色彩噪声和自动曝光特性，渲染出的图像在纹理细节和光照变化上也与真实场景存在差异，这一仿真到真实差距（sim-to-real gap）将在下一章真机部署中进一步讨论。

接触模型方面，本文依赖 MuJoCo 模型自身的接触求解与 Lite3 的开放神经网络交换格式（Open Neural Network Exchange, ONNX）运动策略来形成行走闭环。

\subsection{仿真任务类型}

本文在 MuJoCo 中组织了四类任务：单目标导航、探索模式、状态机全流程导航和多目标任务。单目标导航用于检验基本闭环能力，验证从视觉输入到运动输出的完整链路是否畅通；探索模式用于检验 goal mask 对应的无显式目标条件在部署侧是否可用，即系统能否在没有明确目标图像的情况下生成合理的运动轨迹；状态机导航用于检验从空闲到完成的完整状态流是否连通，验证状态转移逻辑的正确性；多目标任务用于检验导航主机能否持续组织任务队列，验证子目标切换和队列管理的可靠性。

每类任务均记录步数、行驶距离、是否到达目标三项核心指标。
步数反映系统的控制循环次数，行驶距离反映机器人的实际空间位移，到达判定则由预设的距离阈值触发。
此外，所有实验均记录了完整的运行日志，包括每步的推理延迟、轨迹点坐标、速度指令和状态机状态，为后续的结果分析提供了数据基础。

\section{MuJoCo 闭环验证与结果分析}\label{sec:mujoco-results}

\subsection{仿真完整性复核}

在仿真实验前，首先进行 MuJoCo 仿真链路完整性复核。复核目标不是比较算法优劣，而是确认四足机器人、场景、相机、状态机和中层控制接口本身能够形成正确闭环：若机器人无法在 easy/medium/hard 基本场景中物理到达目标，后续算法扫描失败便难以归因于算法本身。

为避免目标判定误触发，在 MuJoCo 任务中显式记录物理目标坐标，并以最终机器人位置到目标点距离不超过 $0.5\,\mathrm{m}$ 作为成功条件；拓扑地图节点预测只在无物理目标的真实部署接口中保留。与此同时，MuJoCo 后端加入仅限仿真域的参考路线稳定项，用于抑制高层航点在走廊墙边产生的横向漂移。该稳定项只在后端为 MuJoCo 仿真且存在物理目标时启用，不改变真机桥接层和真实部署链路。

表~\ref{tab:mujoco-health-check} 给出完整性复核结果。三张地图均满足物理到达阈值，说明后续闭环扫描建立在可运行的 Lite3 MuJoCo 仿真之上。

\begin{table}[htbp]
    \centering
    \caption{Lite3 MuJoCo 完整性复核结果}
    \label{tab:mujoco-health-check}
    \begin{tabular}{lcccc}
        \toprule
        地图 & 步数 & 行驶距离/m & 终点距离/m & 物理到达 \\
        \midrule
        easy & 75 & 5.458 & 0.470 & 是 \\
        medium & 68 & 6.321 & 0.487 & 是 \\
        hard & 110 & 8.442 & 0.475 & 是 \\
        \bottomrule
    \end{tabular}
\end{table}

\subsection{闭环参数扫描实验}

在完整性复核通过后，本文补充了三组仿真实验，实现方式为闭环参数扫描。所有结果均采用 3 次重复实验，并以 $0.5\,\mathrm{m}$ 物理到达阈值判定成功。

第一组实验比较不同 DDIM 步数。表~\ref{tab:mujoco-ddim-closed-loop} 显示，在 medium 与 hard 两个场景中，DDIM-2/3/5/10 均达到 3/3 成功。由于 MuJoCo 参考路线稳定项保证了执行路径的物理可行性，不同步数下的步数和路径长度基本一致；差异主要体现在运行时延上，DDIM-10 明显慢于 DDIM-2。这一结果支持第二章的结论：少步 DDIM 更适合进入部署闭环。

\begin{table}[htbp]
    \centering
    \caption{不同 DDIM 步数的 MuJoCo 闭环实验结果}
    \label{tab:mujoco-ddim-closed-loop}
    \begin{tabular}{lcccc}
        \toprule
        配置 & medium 成功率 & hard 成功率 & medium 时延/s & hard 时延/s \\
        \midrule
        DDIM-2 & 3/3 & 3/3 & 7.5 & 10.0 \\
        DDIM-3 & 3/3 & 3/3 & 7.9 & 10.2 \\
        DDIM-5 & 3/3 & 3/3 & 8.4 & 11.0 \\
        DDIM-10 & 3/3 & 3/3 & 9.9 & 13.5 \\
        \bottomrule
    \end{tabular}
\end{table}

第二组实验固定 DDIM-3，在 medium 场景下比较 CFG 权重。表~\ref{tab:mujoco-cfg-closed-loop} 显示，$w=0,0.5,1.0,2.0$ 均能完成导航，路径指标基本一致；但随着 CFG 权重增大，单次任务运行时间从 7.7\,s 增至 9.2\,s。结合第二章离线结果，CFG 更适合作为可调增强项，而不是默认开启的固定项。

\begin{table}[htbp]
    \centering
    \caption{不同 CFG 权重的 MuJoCo 闭环实验结果（medium, DDIM-3）}
    \label{tab:mujoco-cfg-closed-loop}
    \begin{tabular}{lcccc}
        \toprule
        CFG 权重 & 成功率 & 步数 & 行驶距离/m & 运行时延/s \\
        \midrule
        0.0 & 3/3 & 68.0 & 6.32 & 7.7 \\
        0.5 & 3/3 & 68.0 & 6.32 & 8.7 \\
        1.0 & 3/3 & 68.0 & 6.32 & 9.1 \\
        2.0 & 3/3 & 68.0 & 6.32 & 9.2 \\
        \bottomrule
    \end{tabular}
\end{table}

第三组实验使用四个已训练视觉编码器权重进行 MuJoCo 闭环对比，围绕两个问题展开：其一，四套权重能否通过统一推理模块接入同一闭环系统；其二，接入后的运行时延是否存在明显差异。表~\ref{tab:mujoco-encoder-closed-loop} 显示，四个权重在 easy/medium/hard 场景下均达到 9/9 成功，说明它们均可完成系统级闭环接入。由于参考路线稳定项压缩了路径差异，表中进一步比较运行时延。EfficientNet-B0 的整体时延最低，DINOv2-Small 与 ConvNeXt-Tiny 次之，ResNet-50 在 hard 场景中时延最高。

\begin{table}[htbp]
    \centering
    \caption{跨视觉编码器的 MuJoCo 闭环对比结果}
    \label{tab:mujoco-encoder-closed-loop}
    \begin{tabular}{lcccc}
        \toprule
        编码器 & 总成功率 & easy 时延/s & medium 时延/s & hard 时延/s \\
        \midrule
        EfficientNet-B0 & 9/9 & 7.9 & 7.9 & 10.3 \\
        DINOv2-Small & 9/9 & 9.8 & 10.2 & 12.9 \\
        ConvNeXt-Tiny & 9/9 & 9.8 & 10.1 & 13.4 \\
        ResNet-50 & 9/9 & 10.8 & 11.2 & 15.2 \\
        \bottomrule
    \end{tabular}
\end{table}

\begin{figure}[htbp]
    \centering
    \includegraphics[width=0.8\linewidth]{mujoco.png}
    \caption{仿真实验}
    \label{fig:mujoco}
\end{figure}

\subsection{结果分析}

实验结果说明，第二章的算法层结论已经能够进入系统级闭环，且在仿真完整性复核通过后，DDIM、CFG 与视觉编码器替换均可通过统一推理模块接入 Lite3 MuJoCo 系统。其中，编码器闭环对比进一步回应了前述两个问题：候选权重均可接入同一闭环系统，但不同编码器的运行时延仍会影响后续部署选择。

需要强调的是，MuJoCo 结果只能表明系统级可行性，不能直接替代真机结论。因此本节的结论是：高层 NoMaD、统一推理模块、状态机、导航主机和 Lite3 平台接口已经能够构成稳定可运行的仿真闭环系统，并为下一章真机部署提供可复用接口与参数初值。

\section{本章小结}

本章完成了从算法层到系统层、再到仿真验证层的过渡。通过引入轮式平台与四足平台部署差异、Lite3 平台关键参数、统一推理模块、状态机、导航主机和中层控制映射，本文将第二章得到的动作扩散视觉导航模型整理为可复用的分层闭环系统。该系统将高层策略与平台执行解耦，使 DDIM、CFG、TTS 和视觉编码器等推理配置被封装在策略层内部；状态机通过显式状态定义、恢复与急停入口提高任务可控性；导航主机则通过任务组织与模式切换支持仿真和真机实验。

在此基础上，本章进一步接入 MuJoCo 仿真后端，完成场景设计、接口复核、任务组织和闭环参数扫描。完整性复核表明 easy、medium 和 hard 三类场景中的 Lite3 MuJoCo 链路均能物理到达目标；DDIM、CFG 和视觉编码器闭环扫描则说明第二章筛选出的高层候选配置能够通过统一推理模块进入系统级闭环。MuJoCo 仿真结果不能替代真实平台结论，但它验证了系统接口的可运行性，并为下一章 Jetson Orin 与 Lite3 真机部署提供了可复用接口与参数初始化依据。
